% ==================================================================
% id:       comb-005
% topic:    קומבינטוריקה
% subtopic: בחירה עם חזרות
% solution: full
% ==================================================================
כותבים מספר טבעי כסכום של מספרים טבעיים, כשיש משמעות לסדר המחוברים. למשל, את 3 אפשר לכתוב כך ב-4 דרכים: $3$, $1+2$, $2+1$ ו-$1+1+1$.
\begin{enumerate}
\item בכמה דרכים אפשר לכתוב כך את 6?
\hint{כתבו את 6 כשורה של שש אחדות: $1\ \_\ 1\ \_\ 1\ \_\ 1\ \_\ 1\ \_\ 1$. בין האחדות יש 5 רווחים, ובכל רווח מחליטים אם לשים בו "$+$" או להשאיר אותו ריק.}
\begin{solbox}
נכתוב את 6 כשורה של שש אחדות. בין האחדות יש 5 רווחים:
\[
1\ \_\ 1\ \_\ 1\ \_\ 1\ \_\ 1\ \_\ 1.
\]
בכל רווח מחליטים אם לשים בו "$+$" או להשאיר אותו ריק (ואז שתי האחדות שמשני צדיו שייכות לאותו מחובר). למשל, $1+1\,1+1\,1\,1$ הוא הסכום $1+2+3$. כל בחירה כזו נותנת סכום אחר, וכל סכום מתקבל מבחירה אחת בדיוק. לכל אחד מ-5 הרווחים יש 2 אפשרויות, ולכן לפי עקרון הכפל יש $2^5=\ans{32}$ דרכים.
\end{solbox}
\item בכמה מהן יש בדיוק 3 מחוברים?
\begin{solbox}
סכום עם בדיוק 3 מחוברים מתקבל כששמים "$+$" בבדיוק 2 מתוך 5 הרווחים. בוחרים 2 רווחים מתוך 5, בלי חשיבות לסדר ובלי חזרות, ולכן יש $\binom52=\ans{10}$ דרכים.
\end{solbox}
\end{enumerate}
